\documentclass[11pt]{article}
\usepackage[T1]{fontenc}
\usepackage{lmodern}
\usepackage[margin=0.91in]{geometry}
\usepackage{amsmath,amssymb,amsthm,mathtools,booktabs,array}
\usepackage{xcolor}
\usepackage[colorlinks=true,linkcolor=blue!45!black,urlcolor=blue!45!black,citecolor=blue!45!black]{hyperref}
\usepackage{microtype}
\usepackage{enumitem}
\newtheorem{theorem}{Theorem}[section]
\newtheorem{lemma}[theorem]{Lemma}
\newtheorem{corollary}[theorem]{Corollary}
\newtheorem{proposition}[theorem]{Proposition}
\theoremstyle{remark}\newtheorem{remark}[theorem]{Remark}
\newcommand{\Ind}{\operatorname{Ind}}
\newcommand{\Reg}{\operatorname{Reg}}
\newcommand{\ord}{\operatorname{ord}}
\newcommand{\RH}{\mathrm{RH}}
\newcommand{\Gal}{\operatorname{Gal}}
\newcommand{\one}{\mathbf{1}}
\newcommand{\Q}{\mathbb Q}
\newcommand{\C}{\mathbb C}
\newcommand{\sgn}{\operatorname{sgn}}
\newcommand{\Fix}{\operatorname{Fix}}
\newcommand{\source}{\texttt{nab:q:lowdegree}}
\title{\bfseries Exponentially smaller fields detecting\newline nonabelian Riemann hypotheses\\[6pt]\large A positive induction identity and its iteration in symmetric groups}
\author{Research extension for the ProveIt GRH project}
\date{27 November 2026}
\begin{document}
\maketitle
\begin{abstract}
An open question in the ProveIt GRH article asks whether the subfield of degree
$n!/2$ used to detect all zeros off the critical line in an $S_n$-extension can be
replaced by subfields of smaller degree without invoking Artin holomorphy. We give
an unconditional affirmative answer for $n\ge5$, with one subfield of degree at
most $n!/3^{\lfloor(n-2)/3\rfloor}$. The same argument applies to $A_n$, and a
four-letter refinement improves the degree by another factor of two in one
congruence class. We prove a stronger statement: equality of the zero supports of
the two Dedekind zeta functions at every point other than the pole at $1$, with a
quantitative bound on their multiplicities. The mechanism is an integral positive
induction identity in $A_5$ which can be reused on arbitrarily many disjoint triples
with the same two spare letters. An explicit nonnegative Frobenius weight gives
a corresponding one-sided prime-sum criterion. The arithmetic estimate needed to
establish nonabelian GRH itself is not proved here.
\end{abstract}

\section{Statement and relation to the source}
Let $E/k$ be a finite Galois extension of number fields, $G=\Gal(E/k)$, and
$\pi_U=\Ind_U^G\one$ for $U\le G$. In the case $G=S_n$, the source article
\cite{proveit} proves under Dirichlet GRH that $\RH(\zeta_E)$ is equivalent to
$\RH(\zeta_{E^{\langle(12)\rangle}})$; its question \source{} asks for lower-degree
subfields without Artin holomorphy. The known reflection comparison uses
$2(n-H_n)$, where $H_n=\sum_{j=1}^n1/j$. The source also gives a stronger
multiplicity coefficient $n-1$ for $n\le12$ in that particular reflection
construction. Here we change the detecting subgroup, not the constant for the
transposition subgroup.

\begin{theorem}[Single-field compression]\label{thm:main}
Suppose $G$ contains a copy of $A_n$, $n\ge5$. Write
\[
 r=\left\lfloor\frac{n-2}{3}\right\rfloor,\quad
 \delta=\begin{cases}1& n-3r=4,\\0& n-3r=2\text{ or }3,\end{cases}
 \qquad h=2^\delta3^r,\quad M=3^\delta5^r.
\]
There is an explicitly specified abelian subgroup $H\simeq C_3^r\times C_2^\delta$ of
$G$, with fixed field $K=E^H$, such that
\begin{enumerate}[label=\textup{(\roman*)},nosep]
\item $[K:k]=|G|/h$;
\item $\zeta_E/\zeta_K$ is entire, and $\zeta_K^M/\zeta_E$ is holomorphic
on $\C\setminus\{1\}$ with a pole of exact order $M-1$ at $1$;
\item at every $s_0\ne1$,
\begin{equation}\label{eq:mult}
 0\le \ord_{s_0}\zeta_K\le\ord_{s_0}\zeta_E\le M\ord_{s_0}\zeta_K;
\end{equation}
\item $\RH(\zeta_E)\Longleftrightarrow\RH(\zeta_K)$.
\end{enumerate}
All assertions are unconditional, in particular independent of Dirichlet GRH or
Artin holomorphy.
\end{theorem}

For $G=S_n$, $h\ge3$ and the new index is strictly smaller than $n!/2$; the ratio
of old to new degree is $h/2=\exp((\log3)n/3+O(1))$. For $G=A_n$, the new
index is $n!/(2h)$. The multiplicity bound $M$ is a guaranteed upper bound,
not an asserted optimal value for the iterated construction.

The distinction from the source's abelian theorem is important. The vendored
Lean development exports RH and GRH for Dirichlet $L$-functions. It does not
state GRH for all Dedekind zeta functions or all Artin $L$-functions. The present
proof uses the source's positive-monomial/Heilbronn method (its
\texttt{nab:lem:heilbronn} and \texttt{nab:cor:sn}) and its $A_5$ class computations,
but does not assume or recheck that Lean proof. The new detector is a single
Dedekind zeta function, rather than a signed combination of tests. See
\texttt{SOURCE\_AUDIT.md} for repository paths, inspected blob identifiers,
and scope of the novelty assessment.

\section{The local positive identities}
\begin{lemma}[Five-letter identity]\label{lem:a5}
In $A=A_5$ on $\{1,2,3,4,5\}$, set
\[
 C=\langle(123)\rangle,\quad P=\operatorname{Stab}_A(5)\simeq A_4,\quad
 D=\langle(123),(12)(45)\rangle\simeq S_3.
\]
Let $\epsilon:D\to\{\pm1\}$ be the sign character of its action on
$\{1,2,3\}$. Then $\ker\epsilon=C$, and the following is an equality of virtual
characters of $A$:
\begin{equation}\label{eq:a5}
 5\Ind_C^A\one-\Reg_A=4\Ind_P^A\one+2\Ind_D^A\epsilon.
\end{equation}
\end{lemma}
\begin{proof}
The map $S_3\to D$ extends each odd permutation on the first three letters by
$(45)$, and each even permutation by the identity. This identifies $D$ with
$S_3$ and verifies the kernel assertion. For any $U\le A$ and a character $\lambda$
of $U$, the induced character at $g$ is
\begin{equation}\label{eq:ind}
 (\Ind_U^A\lambda)(g)=\frac{|C_A(g)|}{|U|}
    \sum_{u\in U\cap g^A}\lambda(u).
\end{equation}
The conjugacy classes of $A$ have sizes $1,15,20,12,12$, respectively. The
nonidentity elements of $P$ are three double transpositions and eight
three-cycles; those of $D$ are three double transpositions (sign $-1$) and
 two three-cycles (sign $+1$). Thus the entire calculation is
\[
\begin{array}{c|rrrrr}
\text{class}&1&2^2&3&5_a&5_b\\\hline
|C_A(g)|&60&4&3&5&5\\
\Ind_C^A\one&20&0&2&0&0\\
\Reg_A&60&0&0&0&0\\
\Ind_P^A\one&5&1&2&0&0\\
\Ind_D^A\epsilon&10&-2&1&0&0\\\hline
5\Ind_C^A\one-\Reg_A&40&0&10&0&0\\
4\Ind_P^A\one+2\Ind_D^A\epsilon&40&0&10&0&0
\end{array}
\]
Both split five-cycle classes vanish on all rows. This establishes the identity
on all of $A$ without appealing to the irreducible character table.
\end{proof}

\begin{lemma}[Four-letter refinement]\label{lem:a4}
In $B=A_4$, let $V\simeq C_2^2$ be its normal Klein four subgroup, and let
$J$ be generated by one double transposition. Then
\begin{equation}\label{eq:a4}
 3\Ind_J^B\one-\Reg_B=2\Ind_V^B\one.
\end{equation}
\end{lemma}
\begin{proof}
At the identity the two sides are $3\cdot6-12=6$ and $2\cdot3=6$.
At a double transposition, the centralizer has order $4$, and the induced
characters $\Ind_J^B\one,\Ind_V^B\one$ take values $2,3$. Both sides are $6$.
At each of the two three-cycle classes all the terms vanish.
\end{proof}

\begin{remark}[Optimality of the five-letter step]\label{rem:sharp}
The coefficient $5$ in \eqref{eq:a5} cannot be decreased if the only input is
nonnegative order on all induced linear characters, even if the trivial and
sign character orders are required to vanish. Indeed, inside $S_5$ consider
the $10$-dimensional permutation representation on unordered pairs and its
$5$-dimensional vertex subrepresentation, embedded by
$(x_i)\mapsto(x_i+x_j)_{i<j}$. This map is injective over $\C$ (the equations
$x_i+x_j=0$ for three distinct vertices force all $x_i=0$). Its invariant
orthogonal complement has character $\xi$, degree $5$, and value $-1$ at a
three-cycle: one unordered pair is fixed but two vertices are fixed. Hence
$\langle\xi,\Ind_C^{S_5}\one\rangle=(5-1-1)/3=1$. Both permutation modules
have trivial multiplicity one and sign multiplicity zero, since their point
stabilizers contain odd permutations. The complement therefore has neither
linear constituent. Its restriction to $A_5$ also has no trivial constituent,
by transitivity of $A_5$ on both sets. This is an actual representation, so all
local positivity tests are met. A proposed bound smaller than $5$ would fail
on this abstract order vector. This does not assert existence of an arithmetic
zero realizing that vector, or sharpness of $5^r$ for general $r$.
\end{remark}

\section{Iteration and proof of Theorem \ref{thm:main}}
Choose $r$ disjoint triples $T_1,\ldots,T_r$ in the $n$-letter copy of $A_n$,
and two spare letters $u,v$ outside all triples. Choose a three-cycle $c_i$ on
$T_i$ and let
\[
 C_i=\langle c_i\rangle,\quad H_i=C_1\times\cdots\times C_i,\quad H_0=1.
\]
On $T_i\cup\{u,v\}$ take the $A_5$ subgroup $A_i$ of permutations of those five
letters. Let $P_i$ be the $A_4$ fixing $v$ inside $A_i$, and $D_i$ the copy of
$S_3$ acting on $T_i$ with parity corrected by $(uv)$. Let $\epsilon_i$ be its
sign. Importantly, $A_i$ need not commute with $A_j$ for $i\ne j$; only
$A_i$ and $H_{i-1}$ must commute, which follows from disjoint supports. They
also intersect trivially. Thus $A_iH_{i-1}=A_i\times H_{i-1}$.

Tensor \eqref{eq:a5} by the trivial character of $H_{i-1}$, and induce from
$A_i\times H_{i-1}$ to $G$. Transitivity of induction gives
\begin{equation}\label{eq:iter}
5\pi_{H_i}-\pi_{H_{i-1}}=4\pi_{P_iH_{i-1}}+
 2\Ind_{D_iH_{i-1}}^G(\epsilon_i\boxtimes\one).
\end{equation}
Here $\Reg_{A_i}=\Ind_{1}^{A_i}\one$; this is exactly why the preceding
subfield, not a power of some unrelated zeta function, occurs in the identity.
Define $K_i=E^{H_i}$, $M_i=E^{P_iH_{i-1}}$, and $N_i=E^{D_iH_{i-1}}$.
The last linear character in \eqref{eq:iter} has kernel $H_i$, so $K_i/N_i$
is quadratic. Artin formalism and abelian reciprocity \cite{fgm} give
\begin{align}
 \frac{\zeta_{K_i}(s)^5}{\zeta_{K_{i-1}}(s)}
 &=\zeta_{M_i}(s)^4 L(s,\epsilon_i\boxtimes\one;E/N_i)^2\label{eq:q1}\\
 &=\zeta_{M_i}(s)^4\left(\frac{\zeta_{K_i}(s)}{\zeta_{N_i}(s)}\right)^2.\notag
\end{align}
For clarity, $E/N_i$ may be nonabelian, but the displayed linear character
factors through its quadratic quotient, so its $L$-function is an entire
nontrivial finite-order Hecke function of $N_i$. No assertion about individual
nonmonomial Artin $L$-functions is involved. Equivalently, a pure zeta identity is
\begin{equation}\label{eq:purezeta}
 \zeta_{K_i}^3\zeta_{N_i}^2=\zeta_{K_{i-1}}\zeta_{M_i}^4.
\end{equation}
These identities hold as meromorphic identities on $\C$, including the factors
at ramified primes, by Artin induction formalism. The right-hand side of
\eqref{eq:q1} is holomorphic away from $1$; its pole at $1$ is of order four,
as also follows directly from the simple poles of the quotient on the left.

If $\delta=1$, there are four spare letters, including $u,v$. Let $B$ be the
$A_4$ on these letters and choose $J,V$ as in Lemma \ref{lem:a4}. It commutes
with $H_r$ and intersects it trivially. Tensor \eqref{eq:a4} by $\one_{H_r}$
and induce to $G$:
\[
 3\pi_{JH_r}-\pi_{H_r}=2\pi_{VH_r}.
\]
Put $H=JH_r$ and $W=E^{VH_r}$. This time
\begin{equation}\label{eq:extra}
 \frac{\zeta_{E^H}(s)^3}{\zeta_{K_r}(s)}=\zeta_W(s)^2.
\end{equation}
If $\delta=0$ put $H=H_r$, with no extra step.

For $\delta=0$ multiply \eqref{eq:q1} over $i$ with exponents $5^{i-1}$:
\begin{equation}\label{eq:tel}
 \frac{\zeta_{K_r}^{5^r}}{\zeta_E}=
 \prod_{i=1}^r\left(\zeta_{M_i}^4L(s,\epsilon_i\boxtimes\one;E/N_i)^2\right)^{5^{i-1}}.
\end{equation}
For $\delta=1$, multiply this by the $5^r$-th power of \eqref{eq:extra}.
In both cases the result is the required holomorphic quotient outside $1$,
with a pole of order $M-1$ at $1$ (all Dedekind zeta functions have simple
poles there). Since $E/E^H$ is Galois with abelian group $H$, its zeta quotient
$\zeta_E/\zeta_{E^H}$ is a product of the entire nontrivial Hecke functions
of $E^H$ associated to the characters of $H$. Taking orders at any $s_0\ne1$
now proves \eqref{eq:mult}. Both zeta functions are holomorphic there, and their
zero sets coincide there. Restricting to the open critical strip proves (iv).
The field degree follows from the explicit direct product and the Galois
correspondence. This completes the proof.

\begin{corollary}[Artin nonvanishing away from detected zeros]\label{cor:artin}
For every $s_0\ne1$ such that $\zeta_K(s_0)\ne0$, all irreducible Artin
$L$-functions of $E/k$ are holomorphic and nonzero at $s_0$. In particular
$\RH(\zeta_K)$ implies their zero and pole exclusion in the open critical strip
off the line $\Re s=1/2$.
\end{corollary}
\begin{proof}
By Brauer meromorphic continuation, define the Heilbronn character
\[\Theta=\sum_{\chi\in\operatorname{Irr}(G)}\ord_{s_0}L(s,\chi;E/k)\chi.\]
Its degree is $\ord_{s_0}\zeta_E=0$. Its restriction to any cyclic subgroup
$U$ is a genuine character: Frobenius reciprocity and induction identify
its linear-character multiplicities with orders of Hecke $L$-functions
of $E^U$, all nonnegative at $s_0\ne1$. A genuine character of degree zero
is zero. Thus $\Theta(g)=0$ for all $g$, by taking $U=\langle g\rangle$.
Orthogonality forces every Artin order to be zero. This is the standard
Heilbronn argument \cite{fgm}, with the new input supplied by Theorem
\ref{thm:main}. No claim of holomorphy on the critical line is made.
\end{proof}

\section{The explicit prime statistic}
Assume now $k=\Q$ and $G=S_n$ or $A_n$. Choose a finite set $S$ of rational
primes containing all primes ramified in $E$. For $p\notin S$ let $g_p$ be a
Frobenius conjugacy class in $G$. The number of degree-one primes of $K$
over $p$ equals the permutation character $w(g_p)=\pi_H(g_p)$.

\begin{proposition}[Closed Frobenius weight]\label{prop:weight}
For $G=S_n$ and $H$ as in Theorem \ref{thm:main}, $w(g)=0$ unless $g$ has
cycle type $2^{2e}3^j1^{n-4e-3j}$, for $0\le e\le\delta$ and $0\le j\le r$.
On these classes,
\begin{equation}\label{eq:weight}
 w(g)=\frac{2^{2e}(2e)!\,3^j j!\,(n-4e-3j)!}{2^\delta3^r}
              \binom{r}{j}2^j.
\end{equation}
For $G=A_n$ divide the right-hand side by two. In both cases the weight has
mean one on $G$ and nonnegative integer values.
\end{proposition}
\begin{proof}
There are $\binom{r}{j}2^j$ elements of $H$ of the indicated type (choose
$j$ nonidentity triple rotations, each in two ways, and choose the uniquely
specified $J$ component). The centralizer size in $S_n$ is the numerator
factor before the binomial; \eqref{eq:ind} yields \eqref{eq:weight}. For
$A_n$ the relevant classes do not split: if $e=0$ at least two letters are fixed,
and if $e=1$ one of the disjoint two-cycles is an odd centralizing permutation.
In either case the centralizer has exactly half its $S_n$ size inside $A_n$.
Burnside's lemma yields mean one since the coset action is transitive.
\end{proof}

\begin{corollary}[One-sided prime criterion]\label{cor:prime}
Define
\[
 P(x)=\sum_{\substack{p\le x\\p\notin S}} w(g_p)\log p
             \left(\frac1p-\frac1x\right),\qquad
 R(x)=P(x)-\log x-c,
\]
where $c=\lim_{x\to\infty}(P(x)-\log x)$. Then the following are equivalent:
\begin{enumerate}[label=\textup{(\roman*)},nosep]
\item $\RH(\zeta_E)$;
\item for every $\eta>0$, $R(x)\le C_\eta x^{-1/2+\eta}$ for all sufficiently
large $x$;
\item $R(x)=O_{K,S}(x^{-1/2})$.
\end{enumerate}
More generally, an eventual upper bound $R(x)\le Cx^{\beta-1}$ with
$1/2<\beta<1$ excludes every zero of $\zeta_E$ in $\Re s>\beta$.
\end{corollary}
\begin{proof}
We give the analytic reduction, following \texttt{nab:prop:dedekind} in the source. Write
$\theta(x)=\sum_{p\le x,p\notin S}w(g_p)\log p$. The classical prime ideal
theorem (with its fixed-field error term) gives the existence of $c$ and
\begin{equation}\label{eq:rint}
 R(x)=-\int_x^\infty(\theta(t)-t)t^{-2}\,dt.
\end{equation}
Indeed $P(x)=\int_1^x\theta(t)t^{-2}dt$ and the error integral converges:
$\psi_K(t)-t=O_K(t\exp(-a_K\sqrt{\log t}))$ for large $t$ (a fixed possible
exceptional real zero may be absorbed by enlarging the implicit constant),
and $\psi_K(t)-\theta(t)=O_{K,S}(\sqrt t)$.

For $\Re s>1$ put $D(s)=\sum_{p\notin S}w(g_p)\log p\,p^{-s}$.
The omitted prime powers, higher residue degrees, and the finitely many primes
in $S$ give a holomorphic correction on $\Re s>1/2$; thus
$D(s)=-\zeta_K'(s)/\zeta_K(s)+b(s)$ there, by meromorphic continuation.
Integrating by parts first for $\Re s>1$ gives
\[
 F(s)=\int_1^\infty R(x)x^{-s}\,dx
 =\frac{R(1)}{s-1}+\frac{1}{s-1}\left(\frac{D(s)}s-\frac1{s-1}\right).
\]
At $1$ the apparent double pole cancels by the residue of $D(s)$; the possible
simple pole cancels since $R(x)\to0$, which implies
$\lim_{\sigma\downarrow1}(\sigma-1)F(\sigma)=0$ (split at a sufficiently large
fixed $x$). At a zero $\rho$ of $\zeta_K$ with $\Re\rho>1/2$, the residue of
$F$ is $\ord_\rho\zeta_K/[\rho(1-\rho)]$. It is nonzero and positive for a
real zero in $(1/2,1)$.

Suppose $R(x)\le Cx^{\beta-1}$ eventually. Increase $C$ to make the bound hold
for all $x\ge1$, using continuity on compact intervals. Then the Mellin transform
of $Cx^{\beta-1}-R(x)\ge0$ equals $C/(s-\beta)-F(s)$ for $\Re s>1$.
If its abscissa of convergence were a real $\sigma_0>\beta$, Landau's theorem
for a nonnegative Mellin transform would force a real singularity at $\sigma_0$.
The only possible singularities there are simple poles at real zeros of
$\zeta_K$ in $(\beta,1)$ (classical nonvanishing on $\Re s\ge1$ apart from the
pole). Each has \emph{negative} residue in $C/(s-\beta)-F(s)$, contradicting
nonnegativity on the real axis to the right of the abscissa. Hence the abscissa
is at most $\beta$ (or $-\infty$), so the transform is holomorphic on
$\Re s>\beta$. All zeros of $\zeta_K$ there would be nonremovable poles of $F$,
and are excluded. Theorem \ref{thm:main} transfers the result to $E$.

If $\RH(\zeta_E)$ holds, it holds for $\zeta_K$. The integrated explicit formula
under RH \cite{gm} gives $\int_1^x(\psi_K(t)-t)dt=O_K(x^{3/2})$.
Since the terms not in $\theta$ are $O_{K,S}(\sqrt t)$ by the Chebyshev bound
for prime squares and higher powers, $A(x)=\int_1^x(\theta(t)-t)dt=O_{K,S}(x^{3/2})$.
Integration by parts in \eqref{eq:rint} yields
\[
 R(x)=A(x)x^{-2}-2\int_x^\infty A(t)t^{-3}dt=O_{K,S}(x^{-1/2}).
\]
This proves (i)$\Rightarrow$(iii)$\Rightarrow$(ii); taking
$\beta=1/2+\eta$ arbitrarily close to $1/2$ gives (ii)$\Rightarrow$(i)
by the functional equation. All constants are for fixed fields and $S$.
\end{proof}

For example, when $n=8$ there are two triples and two spare letters, so
$|H|=9$. In an $S_8$-extension the weight is $4480$ at the identity,
$160$ at a single three-cycle, $16$ at a pair of disjoint three-cycles, and
zero otherwise. The corresponding class densities are $1/40320$, $112/40320$,
$1120/40320$; their weighted sum is one. The test does not require that all
primes split completely. In particular, one may take the splitting field of
$T^8-T-1$, whose Galois group is $S_8$ by the factorization and primitivity
certificate in the source (\texttt{grp:ex:s8}). This uses only three
Frobenius patterns, rather than the six in that source example. A bound in (ii) for these explicit weights is still
an open arithmetic problem, not a consequence of the character identity.

\section{The small-degree boundary}
\begin{proposition}[$S_4$: two smaller detectors]\label{prop:s4}
Let $G=S_4$ and $C=\langle(123)\rangle$, $V$ the normal Klein four subgroup.
Then, with $F=E^{A_4}$, $K_8=E^C$, $K_6=E^V$,
\begin{equation}\label{eq:s4}
 \zeta_{K_8}^3\zeta_{K_6}=\zeta_E\zeta_F^3.
\end{equation}
In particular, RH for $\zeta_E$ is equivalent to RH for both $\zeta_{K_8}$
and $\zeta_{K_6}$. The relative degrees $8,6$ are both smaller than $4!/2$.
\end{proposition}
\begin{proof}
The identity of characters $3\pi_C+\pi_V-\Reg_G=3\pi_{A_4}$ follows from
\eqref{eq:ind} on the five cycle types $1^4,2,2^2,3,4$; the left side takes
values $6,0,6,6,0$, respectively. Hence \eqref{eq:s4}. At $s_0\ne1$, if both
smaller zeta orders vanish, the identity forces the order of $\zeta_E$ to vanish
since $\zeta_F$ is holomorphic there. Conversely $E/K_8$ and $E/K_6$ are
abelian, so each smaller order is at most that of $\zeta_E$. The conclusion
follows on the open critical strip. Over $\Q$, Dirichlet GRH implies that
$\ord_\rho\zeta_F=0$ at any point of the strip off the line; in that case
\eqref{eq:s4} gives the exact formula
$\ord_\rho\zeta_E=3\ord_\rho\zeta_{K_8}+\ord_\rho\zeta_{K_6}$.
\end{proof}

\begin{proposition}[Limit of a single-field local positivity method]\label{prop:smallobs}
For $S_3$, no family of subgroups of index strictly smaller than $3$ can
provide zero detection using only monomial-order positivity and vanishing
of the two linear constituents. For $S_4$ no \emph{single} subgroup of index
strictly smaller than $12$ can do so. These are obstructions to that abstract
method, not counterexamples to arithmetic RH.
\end{proposition}
\begin{proof}
In $S_3$, every subgroup of index less than $3$ contains $C_3$. The standard
two-dimensional representation has no invariants under $C_3$, has no linear
constituent, and is genuine, so it supplies the common obstruction.

In $S_4$, a subgroup $U$ of index less than $12$ has order greater than two.
If it contains a three-cycle, use the standard representation of $S_3$ inflated
from $S_4/V\simeq S_3$, which has no invariants under that three-cycle. If it
contains a four-cycle, use the three-dimensional standard representation
of $S_4$ (the sum-zero subspace of its vertex representation). A four-cycle
has no fixed vector there. Otherwise all nonidentity elements of $U$ have
order two, and $U$ is elementary abelian since $(ab)^2=1$ implies $ab=ba$.
If $U$ has no transposition it must be $V$, which acts transitively on vertices,
so the standard representation has no invariants. If it contains a
transposition, say $(12)$, it must lie in the centralizer
$\{1,(12),(34),(12)(34)\}$ and be all of that subgroup. The sign twist of the
standard representation then has no $U$-invariants: the $-1$ eigenspaces
of $(12)$ and $(34)$ in the untwisted representation are two different
coordinate-difference lines with zero intersection. Every representation just
used is genuine, of positive degree, and has neither a trivial nor a sign
constituent. It therefore satisfies the necessary local positivity and
linear-vanishing conditions while its $U$-invariant dimension is zero.
\end{proof}

\section{Quantitative comparison and status}
\begin{center}
\begin{tabular}{rrrrr}
\toprule
$n$&$|H|$&old degree $n!/2$&new degree $n!/|H|$&$M$\\\midrule
5&3&60&40&5\\
6&3&360&240&5\\
7&6&2,520&840&15\\
8&9&20,160&4,480&25\\
10&18&1,814,400&201,600&75\\
12&27&239,500,800&17,740,800&125\\
\bottomrule
\end{tabular}
\end{center}
The algebraic certificate consists of the two small integral identities,
transitivity of induction, and their telescoping product. The included
\texttt{verify.py} independently evaluates the local induced characters on every
group element and checks the class-coefficient identities for every tower step
for $5\le n\le60$ in exact rational arithmetic. The file
\texttt{verification.json} records these checks and the degree formulas. The
proof is symbolic and valid for arbitrary $n$; the finite verification is an
audit, not the basis of the universal quantifier.

\paragraph{What is resolved.} The explicit lower-degree detector question
\source{} of the ProveIt article receives an affirmative answer for all
$n\ge5$ with one field, and for $n=4$ with two fields. The result is stronger
than off-line detection and does not rely on the external Lean theorem. A
particular $A_5$ identity is amplified using shared spare letters; choosing
five disjoint letters for each step would yield only $3^{\lfloor n/5\rfloor}$,
rather than the achieved factor $3^{\lfloor(n-2)/3\rfloor}$. The construction
applies to any ambient finite Galois group containing $A_n$.

\paragraph{What is not resolved.} No new zero-free region for a nonabelian
Dedekind zeta function is claimed, and the distinct open question in the source
about the optimal $n-1$ reflection coefficient in all degrees is not answered.
The article gives a proved reduction of nonabelian RH to a single, smaller
field, or to a nonnegative prime statistic. Establishing the requisite upper
bound for that statistic remains the substantive analytic obstacle. The degree
reduction does not imply an equal reduction of conductor, numerical cost, or
prime-sum difficulty. No claim of global optimality for $H$ or $M$ is made.
The proof and finite certificates have been checked in this research session,
not formally kernel-checked in Lean or independently peer reviewed. The
result is new relative to the inspected source section; an exhaustive priority
search across the mathematical literature has not been performed.

\begin{thebibliography}{9}
\bibitem{proveit} V. Reshetnikov, \emph{ProveIt: Generalized Riemann Hypothesis},
Article, section ``Nonabelian reductions'' (in \texttt{nab.tex}), especially
\texttt{nab:prop:dedekind}, \texttt{nab:lem:heilbronn}, \texttt{nab:cor:sn},
\source{}. Repository main branch inspected 27 November 2026.
\href{https://github.com/VladimirReshetnikov/ProveIt/tree/main/Analysis/GeneralizedRiemannHypothesis}{Repository}.
\bibitem{fgm} R. Foote, H. Ginsberg and V. K. Murty, \emph{Heilbronn characters},
Bull. Amer. Math. Soc. \textbf{52} (2015), 465--496.
\href{https://doi.org/10.1090/bull/1492}{doi:10.1090/bull/1492}.
\bibitem{gm} L. Greni\'e and G. Molteni, \emph{Explicit smoothed prime ideals
 theorems under GRH}, Math. Comp. \textbf{85} (2016), 1875--1899.
\href{https://doi.org/10.1090/mcom3039}{doi:10.1090/mcom3039}.
\end{thebibliography}
\end{document}
